\\documentclass[11pt]{article}
\\usepackage[T1]{fontenc}
\\usepackage[utf8]{inputenc}
\\usepackage{geometry}
\\usepackage{xcolor}
\\usepackage{fvextra}
\\usepackage{hyperref}
\\geometry{margin=1in}
\\DefineVerbatimEnvironment{CodeBlock}{Verbatim}{breaklines,breakanywhere,fontsize=\\small}

\begin{document}
\section*{run\\_experiments.py}
\textbf{Source Path}: \texttt{run\\_experiments.py}\par
\bigskip
\begin{CodeBlock}
"""
Belief Transformer Experimental Runner - COMPLETE VERSION

Supports:
- Multiple experiment modes (standard, enhanced, shared_pca, multi_kernel, etc.)
- Real corpus, temporal batches, AND control corpora (NEW!)
- Multi-kernel experiments  
- Sigma/alpha sweeps
- Batch temporal processing

Usage:
    # Real corpus
    python run_experiments.py --corpus real --mode enhanced --seeds 42 43 44
    
    # Control corpora (NEW!)
    python run_experiments.py --corpus control_constant --mode enhanced --seeds 42 43 44
    python run_experiments.py --corpus control_shuffled --mode enhanced --seeds 42 43 44
    python run_experiments.py --corpus control_random --mode enhanced --seeds 42 43 44
    
    # Temporal batches
    python run_experiments.py --corpus temporal --mode enhanced --seeds 42 43 --batch-temporal
    
    # Legacy control (backward compatible)
    python run_experiments.py --corpus control --mode enhanced --seeds 42 43 44
"""

# ============================================================================
# WINDOWS UNICODE FIX - Must be before all other imports
# ============================================================================
import sys
if sys.platform == 'win32':
    import io
    sys.stdout = io.TextIOWrapper(sys.stdout.buffer, encoding='utf-8', errors='replace')
    sys.stderr = io.TextIOWrapper(sys.stderr.buffer, encoding='utf-8', errors='replace')
# ============================================================================

from pathlib import Path
import json
import argparse, os
from datetime import datetime
import subprocess
import torch
import hashlib

from core.data_utils import extract_article_text
from core.canonical_ids import assign_canonical_uids

# Parsed CLI args (set in main)
args = None

# Project root
ROOT = Path(__file__).resolve().parent
if str(ROOT) not in sys.path:
    sys.path.insert(0, str(ROOT))

DATA_DIR = ROOT / "data"
OUTPUT_DIR = ROOT / "outputs"

from core.complete_pipeline import run_multi_observer_experiment
from core.pipeline_config import PipelineRuntimeConfig


# =========================================================================
# EXPERIMENT CONFIGURATIONS
# =========================================================================

EXPERIMENT_MODES = {
    'standard': {
        'name': 'Standard Mode',
        'description': 'Basic pipeline (no contrastive, no PCA)',
        'use_contrastive': False,
        'use_pca_removal': False,
        'shared_pca': False,
        'kernel_type': 'rbf'
    },
    
    'contrastive': {
        'name': 'Contrastive Only',
        'description': 'Contrastive NLI features only',
        'use_contrastive': True,
        'use_pca_removal': False,
        'shared_pca': False,
        'kernel_type': 'rbf'
    },
    
    'pca': {
        'name': 'PCA Removal Only',
        'description': 'PCA topic removal only',
        'use_contrastive': False,
        'use_pca_removal': True,
        'shared_pca': False,
        'kernel_type': 'rbf'
    },
    
    'enhanced': {
        'name': 'Enhanced Mode',
        'description': 'Contrastive NLI + PCA removal (original)',
        'use_contrastive': True,
        'use_pca_removal': True,
        'shared_pca': False,
        'kernel_type': 'rbf'
    },
    
    'shared_pca': {
        'name': 'Shared PCA Control',
        'description': 'Tests if observer variance is real vs PCA artifact',
        'use_contrastive': True,
        'use_pca_removal': True,
        'shared_pca': True,
        'kernel_type': 'rbf'
    },
    
    'multi_kernel': {
        'name': 'Multi-Kernel Observers',
        'description': 'Different kernels = genuinely different observers',
        'use_contrastive': True,
        'use_pca_removal': True,
        'shared_pca': True,
        'kernel_types': ['rbf', 'laplacian', 'rq', 'imq'],
        'seed': 42
    },
    
    'sigma_sweep': {
        'name': 'Kernel Bandwidth Sweep',
        'description': 'Different sigma values = observers with different scales',
        'use_contrastive': True,
        'use_pca_removal': True,
        'shared_pca': True,
        'kernel_type': 'rbf',
        'sigmas': [0.3, 0.6, 0.9, 1.2, 1.5],
        'seed': 42
    },
    
    'rq_sweep': {
        'name': 'Rational Quadratic alpha Sweep',
        'description': 'Different smoothness parameters',
        'use_contrastive': True,
        'use_pca_removal': True,
        'shared_pca': True,
        'kernel_type': 'rq',
        'alphas': [0.5, 1.0, 2.0, 5.0],
        'sigma': 0.6,
        'seed': 42
    },
    
    'laplacian': {
        'name': 'Laplacian Kernel',
        'description': 'Sharp, local kernel (better for clusters)',
        'use_contrastive': True,
        'use_pca_removal': True,
        'shared_pca': True,
        'kernel_type': 'laplacian'
    },
    
    'framing_kernels': {
        'name': 'Multi-Kernel Per Framing',
        'description': 'Different kernel for each of 8 framings (maintains structure)',
        'use_contrastive': True,
        'use_pca_removal': True,
        'shared_pca': True,
        'use_multi_framing_rks': True,
        'kernel_types_per_framing': ['rbf', 'laplacian', 'rq', 'imq', 'matern', 'rbf', 'laplacian', 'rq']
    },
    
    'framing_kernels_uniform': {
        'name': 'Same Kernel All Framings',
        'description': 'All framings use same kernel (baseline for framing_kernels)',
        'use_contrastive': True,
        'use_pca_removal': True,
        'shared_pca': True,
        'use_multi_framing_rks': True,
        'kernel_type': 'rbf'
    },
    
    'minimal': {
        'name': 'Minimal Architecture',
        'description': 'NLI -> RKS -> Attention (no GRU)',
        'use_contrastive': True,
        'use_pca_removal': True,
        'shared_pca': True,
        'use_gru': False,
        'kernel_type': 'rbf'
    },
    
    'no_gru': {
        'name': 'No Temporal GRU',
        'description': 'Tests if temporal processing affects variance',
        'use_contrastive': True,
        'use_pca_removal': True,
        'shared_pca': True,
        'use_gru': False,
        'kernel_type': 'rbf'
    },
    
    'cls_logits': {
        'name': 'CLS+Logits Stacking',
        'description': 'Full architecture with CLS tokens and logits stacking (8192D)',
        'use_contrastive': True,
        'use_pca_removal': False,
        'shared_pca': False,
        'kernel_type': 'rbf',
        'use_cls_tokens': True,
        'use_attention': True,
        'use_dirichlet_fusion': True,
        'dirichlet_alpha': 1.0,
        'dirichlet_n_observers': 50,
        'projection_dim': 256,
        'apply_pca_to_cls': True,
        'normalize_before_projection': True,
    },
    
    'cls_logits_no_pca': {
        'name': 'CLS+Logits (No PCA)',
        'description': 'CLS stacking without PCA topic removal',
        'use_contrastive': True,
        'use_pca_removal': False,
        'shared_pca': False,
        'kernel_type': 'rbf',
        'use_cls_tokens': True,
        'projection_dim': 256,
        'apply_pca_to_cls': False,
        'normalize_before_projection': True,
    },
    
    'cls_logits_paragraph': {
        'name': 'CLS+Logits Paragraph-Aware',
        'description': 'CLS stacking with paragraph-aware NLI extraction',
        'use_contrastive': True,
        'use_pca_removal': False,
        'shared_pca': False,
        'kernel_type': 'rbf',
        'use_cls_tokens': True,
        'projection_dim': 256,
        'apply_pca_to_cls': True,
        'paragraph_aware': True,
        'paragraph_weights': [0.5, 0.3, 0.2],
        'normalize_before_projection': True,
    },
    
    # =========================================================================
    # NEW MODES (Pipeline Pivot)
    # =========================================================================
    
    'unified': {
        'name': 'Unified Pipeline (Default)',
        'description': 'Single inference: raw logits + CLS Dirichlet fusion + optional head recording',
        'use_unified_pipeline': True,
        'extract_logits': True,
        'extract_cls': True,
        'record_heads': False,  # Set True for diagnostics
        'dirichlet_alpha': 1.0,
        'dirichlet_n_observers': 50,
        'dirichlet_rks_dim': 2048,
    },
    
    'unified_with_heads': {
        'name': 'Unified + Head Recording',
        'description': 'Full unified pipeline with attention head diagnostics',
        'use_unified_pipeline': True,
        'extract_logits': True,
        'extract_cls': True,
        'record_heads': True,
        'dirichlet_alpha': 1.0,
        'dirichlet_n_observers': 50,
        'dirichlet_rks_dim': 2048,
    },
    
    'dirichlet': {
        'name': 'Dirichlet Observer Fusion',
        'description': 'Observers are Dirichlet-weighted mixtures of bot CLS embeddings',
        'use_contrastive': True,
        'use_pca_removal': False,
        'shared_pca': False,
        'kernel_type': 'rbf',
        'use_cls_tokens': True,
        'use_dirichlet_fusion': True,
        'dirichlet_alpha': 1.0,
        'dirichlet_n_observers': 50,
    },
    
    'logits_raw': {
        'name': 'Raw Logits (No Kernel)',
        'description': 'Logits as raw 24D coordinates - no RKS expansion',
        'use_contrastive': True,
        'use_pca_removal': False,
        'shared_pca': False,
        'use_rks': False,  # KEY: no kernel expansion
        'use_gru': False,
        'kernel_type': None,
    },
    
    'cls_dirichlet': {
        'name': 'CLS with Dirichlet Fusion',
        'description': 'CLS embeddings fused via Dirichlet observers (full pipeline)',
        'use_contrastive': True,
        'use_pca_removal': False,
        'shared_pca': False,
        'kernel_type': 'rbf',
        'use_cls_tokens': True,
        'use_dirichlet_fusion': True,
        'dirichlet_alpha': 1.0,
        'dirichlet_n_observers': 50,
        'run_dirichlet_controls': True,
    },
}


# =========================================================================
# HELPER FUNCTIONS
# =========================================================================

def load_articles(filepath):
    """Load articles from JSONL file"""
    if not Path(filepath).exists():
        raise FileNotFoundError(f"Article file not found: {filepath}")
    
    articles = []
    with open(filepath, 'r', encoding='utf-8') as f:
        for line in f:
            line = line.strip()
            if not line:
                continue
            articles.append(json.loads(line))

    # Front-load canonical IDs so every downstream path starts with strict bt_uid keys.
    articles, _uid_stats = assign_canonical_uids(articles)
    return articles


def load_corpus(corpus_type, limit=None):
    """
    Load corpus based on type - UPDATED with path support!
    """
    # Support for custom corpus paths (Direct Injection)
    if os.path.exists(corpus_type):
        articles = load_articles(Path(corpus_type))
        if limit:
            articles = articles[:limit]
        return articles

    if corpus_type == 'control_constant':
        corpus_file = DATA_DIR / 'control_constant.jsonl'
        articles = load_articles(corpus_file)
    elif corpus_type == 'control_shuffled':
        corpus_file = DATA_DIR / 'control_shuffled.jsonl'
        articles = load_articles(corpus_file)
    elif corpus_type == 'control_random':
        corpus_file = DATA_DIR / 'control_random.jsonl'
        articles = load_articles(corpus_file)
    elif corpus_type == 'real':
        corpus_file = DATA_DIR / 'real_corpus.jsonl'
        articles = load_articles(corpus_file)
    elif corpus_type == 'temporal':
        # Temporal is handled separately in main loop usually
        raise ValueError("Temporal corpus should be handled via --batch-temporal")
    else:
        raise ValueError(f"Unknown corpus type: {corpus_type}")

    if limit:
        articles = articles[:limit]
    return articles


def print_experiment_header(mode_name, mode_config):
    """Print experiment header"""
    print("\n" + "="*70)
    print(f"EXPERIMENT: {mode_name}")
    print("="*70)
    print(f"Description: {mode_config['description']}")
    print(f"Contrastive: {mode_config.get('use_contrastive', False)}")
    print(f"PCA Removal: {mode_config.get('use_pca_removal', False)}")
    print(f"Shared PCA: {mode_config.get('shared_pca', False)}")
    print(f"Kernel Type: {mode_config.get('kernel_type', 'rbf')}")
    print("="*70)


def verify_articles(articles, manifest):
    """
    Verify a list of article records against a manifest of expected hashes.

    This helper examines each article's text content, computes an MD5 hash,
    and retains only those whose digest appears in ``manifest['article_hashes']``.
    Articles lacking a textual field are conservatively retained. A summary
    of the verification process is printed to stdout.

    Parameters
    ----------
    articles : list
        Articles loaded from a corpus via ``load_corpus``.
    manifest : dict or None
        Manifest object loaded from JSON. Must contain the key
        ``'article_hashes'`` mapping to a list of hexadecimal digests.

    Returns
    -------
    list
        The filtered list of articles. If no manifest is provided or the
        manifest lacks the ``'article_hashes'`` key, the original
        ``articles`` list is returned unchanged.
    """
    if not manifest or 'article_hashes' not in manifest:
        return articles
    expected_hashes = set(manifest.get('article_hashes', []))
    if not expected_hashes:
        return articles
    verified = []
    discarded = 0
    for article in articles:
        try:
            # Determine the best available textual field for hashing
            if isinstance(article, dict):
                if 'text' in article and isinstance(article['text'], str):
                    body = article['text']
                elif 'content' in article and isinstance(article['content'], str):
                    body = article['content']
                else:
                    # Concatenate all string values as a fallback
                    body_parts = [str(v) for v in article.values() if isinstance(v, str)]
                    body = ''.join(body_parts)
            else:
                # Non-dict articles are retained without verification
                verified.append(article)
                continue
            digest = hashlib.md5(body.encode('utf-8')).hexdigest()
            if digest in expected_hashes:
                verified.append(article)
            else:
                discarded += 1
        except Exception:
            # In case of any unexpected failure, keep the article
            verified.append(article)
    print(f"[OK] Verified articles via manifest: kept {len(verified)} / {len(articles)}")
    if discarded:
        print(f"  Discarded {discarded} articles not present in manifest hashes")
    return verified


def _ensure_verification_provenance(artifact: dict, mode_config: dict, seed: int) -> tuple[dict, bool]:
    """
    Ensure observer artifact contains thesis verifier-required provenance keys:
    basis_hash, crn_seed, alpha, weights_hash.
    """
    if not isinstance(artifact, dict):
        return artifact, False

    meta = artifact.get("meta", {})
    if not isinstance(meta, dict):
        meta = {}
    existing = artifact.get("provenance")
    if not isinstance(existing, dict):
        existing = meta.get("provenance", {})
    if not isinstance(existing, dict):
        existing = {}

    prov = dict(existing)

    # basis_hash fallback from kernel/basis config if not emitted by pipeline
    if not prov.get("basis_hash"):
        basis_src = (
            f"{mode_config.get('kernel_type', 'rbf')}|"
            f"{mode_config.get('dirichlet_rks_dim', mode_config.get('rks_dim', 2048))}|"
            f"{mode_config.get('dirichlet_basis_seed', 42)}"
        )
        prov["basis_hash"] = hashlib.sha256(basis_src.encode("utf-8")).hexdigest()

    if prov.get("crn_seed") is None:
        # Use suite-level CRN default unless caller provided an explicit CRN seed.
        default_crn = getattr(args, "crn_seed", 12345) if args is not None else 12345
        prov["crn_seed"] = int(default_crn)

    if prov.get("alpha") is None:
        prov["alpha"] = float(mode_config.get("dirichlet_alpha", 1.0))

    if not prov.get("weights_hash"):
        # Fallback to a CRN contract fingerprint (invariant across corpora).
        weights_src = (
            f"{prov['basis_hash']}|{prov['alpha']}|{prov['crn_seed']}|"
            f"{mode_config.get('dirichlet_n_observers', 50)}|"
            f"{mode_config.get('kernel_type', 'rbf')}"
        )
        prov["weights_hash"] = hashlib.sha256(weights_src.encode("utf-8")).hexdigest()

    changed = prov != existing
    artifact["provenance"] = prov
    meta["provenance"] = prov
    artifact["meta"] = meta
    return artifact, changed


# =========================================================================
# EXPERIMENT RUNNERS
# =========================================================================

def run_standard_experiment(articles, mode_config, seeds, corpus_name='real', *, output_root: str = None, gru_mode: str = 'intra'):
    """
    Run a standard experiment (single kernel type, multiple seeds).

    Parameters
    ----------
    articles : list
        A list of article records for the experiment.
    mode_config : dict
        Configuration dictionary for the selected experiment mode.
    seeds : list[int]
        List of random seeds to initialise independent observers.
    corpus_name : str, optional
        Name of the corpus being processed. Used in default output paths.
    output_root : str, optional
        If provided, this directory is used as the base output directory for
        saving the results. Each seed will be saved as
        ``Path(output_root) / f'seed_{seed}.pt'``. If ``None`` (default),
        the results will be stored under ``outputs/<corpus_name>/<gru_mode>/``.
    gru_mode : str, optional
        Indicates the GRU aggregation mode (``'intra'`` or ``'inter'``). Used
        only when ``output_root`` is ``None`` to organise the output files.
    """
    print_experiment_header(mode_config['name'], mode_config)
    
    runtime_config = PipelineRuntimeConfig.from_mode_config(mode_config)

    # Prepare config for pipeline
    pipeline_config = {
        'use_contrastive': runtime_config.use_contrastive,
        'use_pca_removal': runtime_config.use_pca_removal,
        'use_cls_tokens': runtime_config.use_cls_tokens,
        'projection_dim': runtime_config.projection_dim,
        'apply_pca_to_cls': runtime_config.apply_pca_to_cls,
        'normalize_before_projection': runtime_config.normalize_before_projection,
        'normalize_features': runtime_config.normalize_features,
        'paragraph_aware': mode_config.get('paragraph_aware', False),
        'paragraph_weights': mode_config.get('paragraph_weights', [0.5, 0.3, 0.2]),
        'kernel_type': runtime_config.kernel_type,
        'kernel_params': {},
        'use_multi_framing_rks': runtime_config.use_multi_framing_rks,
        'use_gru': runtime_config.use_gru,
    }
    
    # Add kernel types per framing if specified
    if 'kernel_types_per_framing' in mode_config:
        pipeline_config['kernel_types'] = mode_config['kernel_types_per_framing']
    
    # Add kernel-specific params
    if mode_config.get('kernel_type') == 'rq':
        pipeline_config['kernel_params']['alpha'] = mode_config.get('alpha', 1.0)
    
    # Run pipeline
    print(f"\nRunning with seeds: {seeds}")
    print(f"Use GRU: {pipeline_config['use_gru']}")
    print(f"Use MultiFramingRKS: {pipeline_config['use_multi_framing_rks']}")
    
    results = run_multi_observer_experiment(
        articles=articles,
        seeds=seeds,
        use_contrastive=pipeline_config['use_contrastive'],
        use_pca_removal=pipeline_config['use_pca_removal'],
        use_cls_tokens=pipeline_config['use_cls_tokens'],
        use_attention=mode_config.get('use_attention', True),
        use_dirichlet_fusion=mode_config.get('use_dirichlet_fusion', False),
        shared_pca=mode_config.get('shared_pca', False),
        kernel_type=pipeline_config.get('kernel_type', 'rbf'),
        kernel_types=pipeline_config.get('kernel_types'),
        kernel_params=pipeline_config.get('kernel_params', {}),
        normalize_features=pipeline_config.get('normalize_features', True),
        use_gru=pipeline_config['use_gru'],
        use_multi_framing_rks=pipeline_config['use_multi_framing_rks'],
        device='cuda' if torch.cuda.is_available() else 'cpu',
        track_variance=args.track_variance,
        output_dir=Path(args.output_root) if args.output_root else Path('outputs'),
        corpus_name=corpus_name,
        nli_cache_path=getattr(args, 'nli_cache_path', None),
        enable_checkpoints=True, # Always capture data lineage
    )
    
    # Save results - BUT only if output_root is NOT provided
    # When output_root IS provided, complete_pipeline.py already saved to observer_{seed}.pt
    mode_slug = mode_config['name'].lower().replace(' ', '_').replace('-', '_')
    output_prefix = f"{corpus_name}_{mode_slug}"
    
    if output_root is None:
        # Legacy mode: save to outputs/ with mode-prefixed names
        for seed in seeds:
            if gru_mode == 'intra' or gru_mode is None:
                output_path = OUTPUT_DIR / f"{output_prefix}_observer_{seed}.pt"
            else:
                base_dir = Path('outputs') / corpus_name / gru_mode
                output_path = base_dir / f"seed_{seed}.pt"
                output_path.parent.mkdir(parents=True, exist_ok=True)

            torch.save(results[seed], output_path)
            print(f"[OK] Saved observer {seed} to {output_path}")
    
    else:
        # When output_root is provided, complete_pipeline already saved files
        for seed in seeds:
            expected_path = Path(output_root) / f"observer_{seed}.pt"
            if expected_path.exists():
                # Backfill verifier-required provenance fields if pipeline omitted them.
                try:
                    artifact = torch.load(expected_path, map_location="cpu", weights_only=False)
                    artifact, changed = _ensure_verification_provenance(artifact, mode_config, seed)
                    if changed:
                        torch.save(artifact, expected_path)
                        print(f"[OK] Added verification provenance: {expected_path.name}")
                except Exception as e:
                    print(f"[WARN] Could not patch provenance for {expected_path.name}: {e}")
                print(f"[OK] Confirmed: {expected_path}")
            else:
                print(f"[WARN] Expected file not found: {expected_path}")
    return results


def run_multi_kernel_experiment(articles, mode_config, corpus_name='real'):
    """Run multi-kernel experiment (different kernels, same seed)"""
    print_experiment_header(mode_config['name'], mode_config)
    
    kernel_types = mode_config['kernel_types']
    seed = mode_config.get('seed', 42)
    
    results_by_kernel = {}
    
    for ktype in kernel_types:
        print(f"\n{'='*70}")
        print(f"KERNEL: {ktype.upper()}")
        print(f"{'='*70}")
        
        pipeline_config = {
            'use_contrastive': mode_config.get('use_contrastive', True),
            'use_pca_removal': mode_config.get('use_pca_removal', True),
            'kernel_type': ktype,
            'kernel_params': {}
        }
        
        if ktype == 'rq':
            pipeline_config['kernel_params']['alpha'] = 1.0
        
        results = run_multi_observer_experiment(
            articles=articles,
            seeds=[seed],
            use_contrastive=pipeline_config['use_contrastive'],
            use_pca_removal=pipeline_config['use_pca_removal'],
            shared_pca=mode_config.get('shared_pca', True),
            kernel_type=pipeline_config['kernel_type'],
            kernel_params=pipeline_config.get('kernel_params', {}),
            device='cuda' if torch.cuda.is_available() else 'cpu',
            track_variance=args.track_variance,
            output_dir=Path(args.output_root) if args.output_root else Path('outputs'),
            corpus_name=corpus_name,
        )
        
        results_by_kernel[ktype] = results
        
        output_path = OUTPUT_DIR / f"{corpus_name}_multikernel_{ktype}_seed{seed}.pt"
        torch.save(results[seed], output_path)
        print(f"[OK] Saved {ktype} observer to {output_path.name}")
    
    return results_by_kernel


def run_sigma_sweep_experiment(articles, mode_config, corpus_name='real'):
    """Run sigma sweep experiment (different bandwidths, same seed)"""
    print_experiment_header(mode_config['name'], mode_config)
    
    sigmas = mode_config['sigmas']
    seed = mode_config.get('seed', 42)
    
    for sigma in sigmas:
        print(f"\n{'='*70}")
        print(f"SIGMA: {sigma}")
        print(f"{'='*70}")
        
        pipeline_config = {
            'use_contrastive': mode_config.get('use_contrastive', True),
            'use_pca_removal': mode_config.get('use_pca_removal', True),
            'kernel_type': mode_config.get('kernel_type', 'rbf'),
            'kernel_params': {},
            'rks_sigma': sigma,
        }
        
        results = run_multi_observer_experiment(
            articles=articles,
            seeds=[seed],
            use_contrastive=pipeline_config['use_contrastive'],
            use_pca_removal=pipeline_config['use_pca_removal'],
            shared_pca=mode_config.get('shared_pca', True),
            kernel_type=pipeline_config['kernel_type'],
            rks_sigma=pipeline_config['rks_sigma'],
            device='cuda' if torch.cuda.is_available() else 'cpu',
            track_variance=args.track_variance,
            output_dir=Path(args.output_root) if args.output_root else Path('outputs'),
            corpus_name=corpus_name,
        )
        
        output_path = OUTPUT_DIR / f"{corpus_name}_sigma{sigma:.1f}_seed{seed}.pt"
        torch.save(results[seed], output_path)
        print(f"[OK] Saved sigma={sigma} observer to {output_path.name}")


def run_rq_sweep_experiment(articles, mode_config, corpus_name='real'):
    """Run RQ alpha sweep experiment (different smoothness, same seed)"""
    print_experiment_header(mode_config['name'], mode_config)
    
    alphas = mode_config['alphas']
    sigma = mode_config.get('sigma', 0.6)
    seed = mode_config.get('seed', 42)
    
    for alpha in alphas:
        print(f"\n{'='*70}")
        print(f"ALPHA: {alpha}")
        print(f"{'='*70}")
        
        pipeline_config = {
            'use_contrastive': mode_config.get('use_contrastive', True),
            'use_pca_removal': mode_config.get('use_pca_removal', True),
            'kernel_type': 'rq',
            'kernel_params': {'alpha': alpha},
            'rks_sigma': sigma,
        }
        
        results = run_multi_observer_experiment(
            articles=articles,
            seeds=[seed],
            use_contrastive=pipeline_config['use_contrastive'],
            use_pca_removal=pipeline_config['use_pca_removal'],
            shared_pca=mode_config.get('shared_pca', True),
            kernel_type='rq',
            kernel_params={'alpha': alpha},
            rks_sigma=pipeline_config['rks_sigma'],
            device='cuda' if torch.cuda.is_available() else 'cpu',
            track_variance=args.track_variance,
            output_dir=Path(args.output_root) if args.output_root else Path('outputs'),
            corpus_name=corpus_name,
        )
        
        output_path = OUTPUT_DIR / f"{corpus_name}_rq_alpha{alpha:.1f}_seed{seed}.pt"
        torch.save(results[seed], output_path)
        print(f"[OK] Saved alpha={alpha} observer to {output_path.name}")


def run_dirichlet_fusion_experiment(articles, mode_config, seeds, corpus_name='real', *, output_root: str = None):
    """
    Run Dirichlet observer fusion experiment.
    
    This implements the new observer model where observers are Dirichlet-weighted
    mixtures of bot CLS embeddings, not random RKS projections.
    """
    print_experiment_header(mode_config['name'], mode_config)
    
    # Import dirichlet fusion module
    try:
        from core.dirichlet_fusion import DirichletFusionConfig, DirichletFusion
    except ImportError as e:
        print(f"ERROR: Failed to import dirichlet_fusion: {e}")
        print("Make sure dirichlet_fusion.py is in core/")
        return {}
    
    # Import NLI extractor
    from core.nli_extraction import UnifiedNLIExtractor, ExtractionConfig
    
    output_dir = Path(output_root) if output_root else OUTPUT_DIR
    output_dir.mkdir(parents=True, exist_ok=True)
    
    device = 'cuda' if torch.cuda.is_available() else 'cpu'
    
    # Create NLI extractor configured for CLS tokens
    extract_config = ExtractionConfig(
        use_cls_tokens=True,  # Correct field name
        extract_logits=False,
    )
    nli_extractor = UnifiedNLIExtractor(config=extract_config)
    
    # Extract CLS embeddings
    print(f"[NLI] Extracting CLS embeddings from {len(articles)} articles...")
    print(f"[NLI] Hidden size: {nli_extractor.hidden_size}")
    cls_per_bot_list = []
    
    batch_size = 512
    for i in range(0, len(articles), batch_size):
        batch = articles[i:i+batch_size]
        texts = [a.get('body', a.get('text', a.get('content', '')))[:2000] for a in batch]
        
        with torch.no_grad():
            result = nli_extractor.extract_batch(texts)
            # Check various possible keys
            if 'cls_per_bot' in result:
                cls_per_bot_list.append(result['cls_per_bot'])
            elif 'cls_views' in result:
                cls_per_bot_list.append(result['cls_views'])
            elif 'cls' in result:
                cls_per_bot_list.append(result['cls'])
        
        if (i // batch_size) % 50 == 0:
            print(f"  Processed {i+len(batch)}/{len(articles)} articles")
    
    cls_per_bot = torch.cat(cls_per_bot_list, dim=0).to(device)  # [N, 8, H]
    actual_hidden = cls_per_bot.shape[-1]
    print(f"[NLI] CLS embeddings shape: {cls_per_bot.shape}")
    

    # Handle ablations
    alpha_val = mode_config.get('dirichlet_alpha', 1.0)
    if getattr(args, 'alpha_collapse', False):
        print("[ABLATION] FORCING ALPHA COLLAPSE (alpha=1e6)")
        alpha_val = 1e6
        
    crn_enabled = True
    if getattr(args, 'no_crn', False):
        print("[ABLATION] DISABLING CRN")
        crn_enabled = False

    # Now configure fusion with actual hidden dim
    config = DirichletFusionConfig(
        n_bots=8,
        hidden_dim=actual_hidden,
        rks_dim=mode_config.get('dirichlet_rks_dim', 2048),
        n_observers=mode_config.get('dirichlet_n_observers', 50),
        alpha=alpha_val,
        kernel_type=mode_config.get('kernel_type', 'rbf'),
        basis_seed=mode_config.get('dirichlet_basis_seed', 42),
        crn_enabled=crn_enabled,
    )
    
    # Create fusion module and run
    fusion = DirichletFusion(config)
    
    print(f"[FUSION] Running Dirichlet fusion with alpha={config.alpha}, K={config.n_observers}...")
    result = fusion(cls_per_bot, return_samples=True, compute_curvature=True)
    
    # Save results
    for seed in seeds:
        output_artifact = {
            'embeddings': result['fused'].cpu(),
            'features': result['fused'].cpu(),
            'observer_samples': result.get('observer_samples', torch.tensor([])).cpu(),
            'fused_std': result['fused_std'].cpu(),
            'bot_rkhs': result.get('bot_rkhs', result.get('rkhs_views', torch.tensor([]))).cpu(),
            'curvature': result.get('curvature', {}),
            'seed': seed,
            'n_articles': result['fused'].shape[0],
            'meta': {
                'mode': 'dirichlet',
                'corpus': corpus_name,
                'config': {
                    'alpha': config.alpha,
                    'n_observers': config.n_observers,
                    'rks_dim': config.rks_dim,
                    'kernel_type': config.kernel_type,
                },
                'provenance': result.get('provenance', {}),
            },
        }
        
        # --- NEW: Preserve full physics payload for bridge extraction ---
        physics_keys = [
            'walker_states', 'walker_work_integrals', 
            'spectral_evr', 'spectral_u_axis', 
            'article_metadata', 'phantom_verdicts'
        ]
        for k in physics_keys:
            if k in result:
                val = result[k]
                # Convert tensors to CPU if needed
                if torch.is_tensor(val):
                    output_artifact[k] = val.cpu()
                else:
                    output_artifact[k] = val
        
        output_artifact, _ = _ensure_verification_provenance(output_artifact, mode_config, seed)
        
        output_file = output_dir / f"observer_{seed}.pt"
        torch.save(output_artifact, output_file)
        print(f"[OK] Saved Dirichlet observer to {output_file}")
    
    return result


def run_unified_pipeline_experiment(articles, mode_config, seeds, corpus_name='real', *, output_root: str = None):
    """
    Run unified pipeline: single inference for logits + CLS Dirichlet + optional heads.
    
    This is the new default mode implementing the pipeline pivot.
    """
    print_experiment_header(mode_config['name'], mode_config)
    
    try:
        from core.unified_extraction import (
            UnifiedExtractionConfig,
            run_unified_pipeline,
        )
    except ImportError:
        print("ERROR: unified_extraction.py not found in core/")
        print("Copy it from the provided files.")
        return {}
    
    output_dir = Path(output_root) if output_root else OUTPUT_DIR
    output_dir.mkdir(parents=True, exist_ok=True)
    
    # Configure
    config = UnifiedExtractionConfig(
        extract_logits=mode_config.get('extract_logits', True),
        extract_cls=mode_config.get('extract_cls', True),
        record_heads=mode_config.get('record_heads', False),
        paragraph_aware=mode_config.get('paragraph_aware', True),
        dirichlet_alpha=mode_config.get('dirichlet_alpha', 1.0),
        dirichlet_n_observers=mode_config.get('dirichlet_n_observers', 50),
        dirichlet_rks_dim=mode_config.get('dirichlet_rks_dim', 2048),
    )
    
    # Run for each seed
    results = {}
    for seed in seeds:
        print(f"\n{'='*60}")
        print(f"SEED: {seed}")
        print(f"{'='*60}")
        
        config.dirichlet_seed = seed
        
        result = run_unified_pipeline(
            articles=articles,
            config=config,
            output_dir=output_dir,
            seed=seed,
        )
        results[seed] = result
    
    return results


def _git_head_short() -> str:
    try:
        proc = subprocess.run(
            ["git", "rev-parse", "--short", "HEAD"],
            check=True,
            capture_output=True,
            text=True,
        )
        return proc.stdout.strip()
    except Exception:
        return "unknown"


def _maybe_freeze_known_good_tuple(
    *,
    output_root: str | None,
    corpus_name: str,
    mode_name: str,
    argv: list[str],
    require_success: bool,
) -> None:
    """Freeze/verify a run tuple when MONOLITH HTML output is present."""
    if output_root:
        run_dir = Path(output_root)
    else:
        corpus_dir = OUTPUT_DIR / corpus_name
        run_dir = corpus_dir if corpus_dir.exists() else OUTPUT_DIR

    if not run_dir.exists():
        if require_success:
            raise RuntimeError(f"[FREEZE] Output directory not found: {run_dir}")
        print(f"[FREEZE] Skip: output directory not found: {run_dir}")
        return

    html_candidates = sorted(run_dir.glob("MONOLITH*.html"), key=lambda p: p.stat().st_mtime, reverse=True)
    if not html_candidates:
        if require_success:
            raise RuntimeError(f"[FREEZE] No MONOLITH HTML found in: {run_dir}")
        print(f"[FREEZE] Skip: no MONOLITH HTML found in {run_dir}")
        return

    source_html = html_candidates[0]
    source_log = source_html.with_suffix(".log")
    if not source_log.exists():
        source_log.write_text("[FREEZE] auto-generated placeholder log\n", encoding="utf-8")

    try:
        from analysis.freeze_viz_tuple import FreezeConfig, freeze_run_snapshot, verify_snapshot
    except Exception as e:
        if require_success:
            raise RuntimeError(f"[FREEZE] Failed to import freeze tooling: {e}") from e
        print(f"[FREEZE] Skip: freeze tooling unavailable ({e})")
        return

    viz_commit = _git_head_short()
    timestamp = datetime.utcnow().strftime("%Y-%m-%d_%H%M%S")
    snapshot_name = f"{timestamp}_{viz_commit}_{corpus_name}_{mode_name}".replace("/", "_").replace(" ", "_")
    command = "python " + " ".join(argv)

    snap_dir = freeze_run_snapshot(
        FreezeConfig(
            snapshot_name=snapshot_name,
            snapshot_root=Path("analysis/locked_runs"),
            source_html=source_html,
            source_log=source_log,
            source_viz_code=Path("analysis/MONOLITH_VIZ.py"),
            source_viz_code_git=Path("analysis/MONOLITH_VIZ.py"),
            input_dir=run_dir,
            viz_code_commit=viz_commit,
            command=command,
        )
    )
    manifest_path = snap_dir / "RUN_MANIFEST.json"
    verify_snapshot(manifest_path)
    print(f"[FREEZE] Snapshot written+verified: {manifest_path}")


# =========================================================================
# MAIN
# =========================================================================

def main():
    parser = argparse.ArgumentParser(
        description="Belief Transformer Experimental Runner"
    )
    
    parser.add_argument(
        '--corpus',
        type=str,
        default='real',
        help='Which corpus to use (NEW: separate control types!)'
    )
    
    parser.add_argument(
        '--mode',
        type=str,
        default='enhanced',
        choices=list(EXPERIMENT_MODES.keys()),
        help='Experiment mode'
    )
    
    parser.add_argument(
        '--seeds',
        type=int,
        nargs='+',
        default=[42, 43, 44, 45, 46],
        help='Observer seeds (for standard modes)'
    )
    
    parser.add_argument(
        '--use-contrastive',
        action='store_true',
        help='Force contrastive NLI (overrides mode)'
    )
    
    parser.add_argument(
        '--use-pca-removal',
        action='store_true',
        help='Force PCA removal (overrides mode)'
    )
    
    parser.add_argument(
        '--queries-config',
        type=str,
        default=None,
        help='Path to contrastive queries config'
    )
    
    parser.add_argument(
        '--compare-modes',
        action='store_true',
        help='Run all modes and compare'
    )
    
    parser.add_argument(
        '--limit',
        type=int,
        default=None,
        help='Limit number of articles (for testing)'
    )
    
    parser.add_argument(
        '--batch-temporal',
        action='store_true',
        help='Process temporal batches separately'
    )
    
    parser.add_argument(
        '--track-variance',
        action='store_true',
        help='Enable variance tracking at each pipeline stage'
    )

    # ------------------------------------------------------------------
    # NEW ARGUMENTS
    #
    # --gru-mode: Choose GRU inter/intra attention mode. Defaults to intra.
    # --output-root: Optional root for structured experiment outputs.
    # --manifest: Optional path to a manifest JSON file containing
    #             article hashes for verifying control corpora.
    #
    # These arguments are added to support more flexible experiment
    # configuration without breaking backwards compatibility. If
    # unspecified, the runner behaves exactly as before.
    parser.add_argument(
        '--gru-mode',
        type=str,
        choices=['intra', 'inter'],
        default='intra',
        help='GRU aggregation mode (intra or inter)'
    )
    parser.add_argument(
        '--output-root',
        type=str,
        default=None,
        help='Root directory for structured output files'
    )
    parser.add_argument(
        '--manifest',
        type=str,
        default=None,
        help='Path to manifest file (for control corpus verification)'
    )
    parser.add_argument(
        '--kernel-type',
        type=str,
        choices=['rbf', 'laplacian', 'rq', 'imq', 'matern', 'cosine', 'linear'],
        default=None,
        help='Override kernel type (default: use mode default)'
    )
    parser.add_argument(
        '--bandwidth',
        type=float,
        default=None,
        help='Kernel bandwidth (sigma) for RBF/Laplacian kernels'
    )
    parser.add_argument(
        '--rks-dim',
        type=int,
        default=2048,
        help='RKS output dimension (default: 2048)'
    )
    parser.add_argument(
        '--kernel-seed',
        type=int,
        default=42,
        help='Seed for RKS basis generation (M-observer control)'
    )
    parser.add_argument(
        '--mix-in-rkhs',
        action='store_true',
        help='Mode B: Map to RKHS before Dirichlet mixing (born-aligned)'
    )
    parser.add_argument(
        '--nli-cache-path',
        type=str,
        default=None,
        help='Path to NLI embedding cache file (for reuse across kernel runs)'
    )
    parser.add_argument(
        '--nu',
        type=float,
        default=None,
        help='Smoothness parameter for Matern kernel (v)'
    )
    parser.add_argument(
        '--roughness',
        type=float,
        default=None,
        help='Roughness parameter for IMQ kernel (df)'
    )
    parser.add_argument(
        '--freeze-good-run',
        action='store_true',
        default=True,
        help='Auto-freeze and verify MONOLITH run tuple after successful run (default: on).'
    )
    parser.add_argument(
        '--no-freeze-good-run',
        action='store_false',
        dest='freeze_good_run',
        help='Disable post-run tuple freeze.'
    )
    parser.add_argument(
        '--freeze-require-success',
        action='store_true',
        help='Fail run if post-run tuple freeze cannot be completed.'
    )
    
    # Dirichlet fusion flags
    parser.add_argument(
        '--dirichlet-fusion',
        action='store_true',
        help='Enable Dirichlet fusion for CLS views'
    )
    parser.add_argument(
        '--alpha',
        type=float,
        default=1.0,
        help='Dirichlet alpha parameter (default: 1.0)'
    )
    parser.add_argument(
        '--fusion-samples',
        type=int,
        default=50,
        help='Number of Dirichlet observer samples (default: 50)'
    )
    parser.add_argument(
        '--shared-basis-path',
        type=str,
        default=None,
        help='Path to shared RKS basis file for CRN reproducibility'
    )
    parser.add_argument(
        '--locked-weights-path',
        type=str,
        default=None,
        help='Path to locked Dirichlet weights file for CRN reproducibility'
    )

    parser.add_argument(
        "--no-crn",
        action="store_true",
        help="Disable CRN for Dirichlet fusion"
    )

    parser.add_argument(
        "--alpha-collapse",
        action="store_true",
        help="Force alpha=1e6 for collapse ablation"
    )
    
    global args
    args = parser.parse_args()
    
    # Resolve corpus name for custom paths
    corpus_name = Path(args.corpus).stem if os.path.exists(args.corpus) else args.corpus


    # ------------------------------------------------------------------
    # Load manifest (if provided)
    #
    # If a manifest JSON file is supplied via --manifest, we read it
    # here and optionally use it later to verify that the articles in a
    # control corpus match the expected set of hashes. This is useful for
    # reproducibility and to detect any accidental drift in the control
    # datasets. The manifest is assumed to contain a key called
    # 'article_hashes' which is a list of hex-digests. If this key is
    # missing, the manifest is ignored.
    manifest = None
    if args.manifest:
        manifest_path = Path(args.manifest)
        if not manifest_path.exists():
            raise FileNotFoundError(f"Manifest file not found: {manifest_path}")
        with open(manifest_path, 'r', encoding='utf-8') as f:
            try:
                manifest = json.load(f)
            except Exception as e:
                raise ValueError(f"Failed to parse manifest JSON: {e}")

    
    # Print header
    print("="*70)
    print("BELIEF TRANSFORMER EXPERIMENTAL RUNNER")
    print("="*70)
    print(f"Corpus: {args.corpus}")
    print(f"Mode: {args.mode}")
    print(f"Seeds: {args.seeds}")
    if args.limit:
        print(f"Limit: {args.limit} articles")
    print("="*70)
    import sys; sys.stdout.flush()  # Ensure output visible in subprocess
    
    # Get mode config
    mode_config = EXPERIMENT_MODES[args.mode].copy()
    
    # Apply overrides
    if args.use_contrastive:
        mode_config['use_contrastive'] = True
    if args.use_pca_removal:
        mode_config['use_pca_removal'] = True
    
    # Kernel type override (for matrix experiments)
    if args.kernel_type:
        mode_config['kernel_type'] = args.kernel_type
        # Clear any multi-kernel settings since we're forcing a specific kernel
        mode_config.pop('kernel_types', None)
        mode_config.pop('kernel_types_per_framing', None)
        print(f"Kernel override: {args.kernel_type}")
    
    # NEW: Kernel configuration for KernelContext
    if args.bandwidth is not None:
        mode_config['kernel_bandwidth'] = args.bandwidth
        print(f"[CONFIG] Kernel bandwidth: {args.bandwidth}")
    
    if args.rks_dim:
        mode_config['dirichlet_rks_dim'] = args.rks_dim
        print(f"[CONFIG] RKS dimension: {args.rks_dim}")
    
    if args.kernel_seed != 42:  # Only print if non-default
        mode_config['dirichlet_basis_seed'] = args.kernel_seed
        print(f"[CONFIG] Kernel seed: {args.kernel_seed}")
    
    # NEW: Mix-in-RKHS (Mode B) - now default for Dirichlet
    if args.mix_in_rkhs:
        mode_config['mix_in_rkhs'] = True
        print("[CONFIG] Mode B enabled: map-then-mix (born-aligned RKHS)")
    
    # NEW: Dirichlet fusion config
    if args.dirichlet_fusion:
        if not mode_config.get('use_cls_tokens', False):
            print("[WARNING] Dirichlet fusion requires CLS tokens. Enabling CLS mode.")
            mode_config['use_cls_tokens'] = True
        
        mode_config['use_dirichlet_fusion'] = True
        mode_config['dirichlet_alpha'] = args.alpha
        mode_config['dirichlet_n_observers'] = args.fusion_samples
        mode_config['mix_in_rkhs'] = True  # Mode B is the only mode now
        
        if args.shared_basis_path:
            mode_config['dirichlet_basis_path'] = args.shared_basis_path
        if args.locked_weights_path:
            mode_config['dirichlet_weights_path'] = args.locked_weights_path
        
        print(f"[CONFIG] Dirichlet fusion enabled: alpha={args.alpha}, K={args.fusion_samples}")
    
    # BATCH TEMPORAL PROCESSING
    if args.batch_temporal and args.corpus == 'temporal':
        print("\n" + "="*70)
        print("[INFO] BATCH PROCESSING MODE - Temporal Evolution")
        print("="*70)
        
        temporal_dir = DATA_DIR / 'temporal_cleaned'
        
        if not temporal_dir.exists():
            raise FileNotFoundError(f"Temporal directory not found: {temporal_dir}")
        
        batch_files = sorted(temporal_dir.glob('*.jsonl'))
        
        if not batch_files:
            raise FileNotFoundError(f"No batch files found in {temporal_dir}")
        
        print(f"Found {len(batch_files)} temporal batches\n")
        
        for i, batch_file in enumerate(batch_files):
            # Extract date range from filename
            date_range = batch_file.stem
            
            print("="*70)
            print(f"BATCH {i+1}/{len(batch_files)}: {date_range}")
            print("="*70)
            
            # Load batch
            batch_articles = load_articles(batch_file)
            # Verify batch via manifest if provided
            if manifest:
                batch_articles = verify_articles(batch_articles, manifest)
            
            if args.limit:
                batch_articles = batch_articles[:args.limit]
            
            print(f"Articles: {len(batch_articles)}")
            
            # Run experiment
            corpus_name = f"temporal_{date_range}"
            
            if 'kernel_types' in mode_config:
                run_multi_kernel_experiment(batch_articles, mode_config, corpus_name)
            elif 'sigmas' in mode_config:
                run_sigma_sweep_experiment(batch_articles, mode_config, corpus_name)
            elif 'alphas' in mode_config:
                run_rq_sweep_experiment(batch_articles, mode_config, corpus_name)
            else:
                run_standard_experiment(
                    batch_articles,
                    mode_config,
                    args.seeds,
                    corpus_name,
                    output_root=args.output_root,
                    gru_mode=args.gru_mode
                )
            
            print(f"\n[OK] Batch {date_range} complete\n")
        
        print("="*70)
        print("[OK] ALL TEMPORAL BATCHES COMPLETE")
        print("="*70)
        return
    
    # NORMAL PROCESSING
    articles = load_corpus(args.corpus, limit=args.limit)
    # Verify articles via manifest if provided
    if manifest:
        articles = verify_articles(articles, manifest)

    # Run appropriate experiment (single-channel only).
    # NOTE: This runner executes exactly ONE mode per invocation.
    #       CLS vs logits separation is handled by the outer suite/orchestrator.
    # Policy:
    #   - NEVER apply pipeline PCA removal to logits modes (non-CLS).
    #   - CLS-only PCA (if enabled) may occur inside the CLS extractor/config.
    if not mode_config.get("use_cls_tokens", False):
        # Hard rule: logits/24D modes never use pipeline PCA removal
        mode_config["use_pca_removal"] = False

    # Optional kernel override (forces single-kernel run even if the mode defines a sweep)
    if getattr(args, "kernel_type", None):
        mode_config["kernel_type"] = args.kernel_type
        mode_config.pop("kernel_types", None)
        mode_config.pop("kernel_types_per_framing", None)
        if 'pipeline_config' in locals():
            pipeline_config['kernel_type'] = args.kernel_type
        print(f"Kernel override: {args.kernel_type}")

    # NEW: Handle unified pipeline mode (single inference for all outputs)
    if mode_config.get('use_unified_pipeline', False):
        run_unified_pipeline_experiment(
            articles,
            mode_config,
            args.seeds,
            corpus_name=corpus_name,
            output_root=args.output_root,
        )
    # Handle Dirichlet fusion standalone mode (only if strictly specified as the mode)
    elif args.mode == 'dirichlet':
        run_dirichlet_fusion_experiment(
            articles,
            mode_config,
            args.seeds,
            corpus_name=corpus_name,
            output_root=args.output_root,
        )
    elif 'kernel_types' in mode_config:
        # Mode-defined multi-kernel sweep (legacy)
        run_multi_kernel_experiment(articles, mode_config, args.corpus)
    elif 'sigmas' in mode_config:
        run_sigma_sweep_experiment(articles, mode_config, args.corpus)
    elif 'alphas' in mode_config:
        run_rq_sweep_experiment(articles, mode_config, args.corpus)
    else:
        run_standard_experiment(
            articles,
            mode_config,
            args.seeds,
            corpus_name=corpus_name,
            output_root=args.output_root,
            gru_mode=getattr(args, "gru_mode", "intra")
        )

    print(f"\nResults saved to: {OUTPUT_DIR}")

    if args.freeze_good_run:
        _maybe_freeze_known_good_tuple(
            output_root=args.output_root,
            corpus_name=corpus_name,
            mode_name=args.mode,
            argv=sys.argv,
            require_success=bool(args.freeze_require_success),
        )
    
    # Control-specific next steps
    if args.corpus.startswith('control_'):
        print(f"\nNext steps:")
        print(f"  1. Run other controls if not done:")
        print(f"     python run_experiments.py --corpus control_constant --mode {args.mode} --seeds {' '.join(map(str, args.seeds))}")
        print(f"     python run_experiments.py --corpus control_shuffled --mode {args.mode} --seeds {' '.join(map(str, args.seeds))}")
        print(f"     python run_experiments.py --corpus control_random --mode {args.mode} --seeds {' '.join(map(str, args.seeds))}")
        print(f"  2. Compare results:")
        print(f"     python controls/compare_controls.py --mode {args.mode}")
    else:
        print(f"\nNext steps:")
        print(f"  1. Run Procrustes alignment:")
        print(f"     cd analysis && python procrustes_alignment.py")
        print(f"  2. Generate visualization:")
        print(f"     python rks_viz_CLEAN.py")
    
    print("="*70)


if __name__ == "__main__":
    main()

\end{CodeBlock}
\end{document}
